% Implementation: imitation/algos/dice/distill_rl.py::field_actor_loss,
% field_pointwise with unit_normalize=false and adaptive_bc=false.
% q_max_norm clips BEFORE q_step_size; total_max_norm clips the combined
% displacement. Both bc_step_size and restore_step_size restore to the base.

\method{} improves sampled actions and trains the residual to reproduce
them. Each target combines a critic-guided step with a restoring correction
toward the frozen base. This gives the update two separate controls:
clipping limits the proposed step, while restoration opposes accumulated
departure. All operations below use normalized action coordinates, with
chunk dimension $D=Hd$.

\subsection{Critic-Guided Proposals}
\label{sec:reward-field}

For each replay state $s$, draw $K$ noise samples and form current actions:
\begin{equation}
  \bar r_i=\operatorname{sg}[r_\theta(s,z_i)],
  \qquad a_i=\pi_0(s,z_i)+\bar r_i,
  \label{eq:candidates}
\end{equation}
where $\operatorname{sg}$ denotes stop-gradient. Let $\bar Q_\phi(s,z,a)$
be the conservative critic-ensemble aggregate; the noise input is omitted
when the critic does not use it. Following the baseline's value
normalization~\cite{dicerl2026}, divide by a detached scale
$q_{\mathrm{scale}}=\max(\epsilon,\mathbb{E}_{\mathrm{batch}}[|\bar Q_\phi|])$,
computed over the current candidates. The reward displacement is
\begin{equation}
  d_i^Q=\eta_Q\,C_{\delta_Q}\!\left(
    \frac{\nabla_a\bar Q_\phi(s,z_i,a_i)}{q_{\mathrm{scale}}}
  \right),
  \label{eq:reward-step}
\end{equation}
where $C_\delta(v)=v\min(1,\delta/\|v\|_2)$, with $C_\delta(0)=0$.
The clip rescales large gradients while preserving small ones; $\eta_Q$
then sets the reward step size. We differentiate the critic with respect
to its action input only, holding its parameters fixed during target
construction.

\subsection{Restoration Toward the Frozen Base}
\label{sec:anchor}

Each candidate is anchored to the base action generated by the same
$(s,z_i)$. Define its residual magnitude in root-mean-square units as
$m_i=\|\bar r_i\|_2/\sqrt{D}$. The restoring displacement is
\begin{equation}
  d_i^R=-\eta_0\bar r_i
        -\eta_{\mathrm{anc}}
        \left(1-\frac{\rho}{\max(m_i,\epsilon)}\right)_{+}\bar r_i,
  \label{eq:restoring-step}
\end{equation}
where $(x)_+=\max(x,0)$. The first term supplies a weak, always-on pull.
The second is inactive within radius $\rho$ and opposes the excess
displacement outside it. Measuring $\rho$ in rms units makes it comparable
across action-chunk dimensions. The reference remains frozen throughout
training, so the restoring force responds to total departure rather than
just the latest step.

This is a soft correction evaluated at the current residual, rather than
the hard projection used for illustration in Fig.~\ref{fig:problem}.
Setting $\rho=0$ makes both terms always active. Removing restoration
requires setting \emph{both} $\eta_0$ and $\eta_{\mathrm{anc}}$ to zero.

\subsection{Regression onto the Adjusted Targets}
\label{sec:field-update}

Combine the two displacements, clip their sum, and detach the target:
\begin{equation}
  \tilde r_i=\operatorname{sg}\!\left[
    \bar r_i+C_{\delta_{\mathrm{tot}}}(d_i^Q+d_i^R)
  \right].
  \label{eq:field-target}
\end{equation}
The residual network is then trained by mean-squared error,
\begin{equation}
  \mathcal L_{\mathrm{actor}}
    =\mathbb{E}_{s}\!\left[
      \frac{1}{KD}\sum_{i=1}^{K}
      \|r_\theta(s,z_i)-\tilde r_i\|_2^2
    \right].
  \label{eq:actor-regression}
\end{equation}
Only this regression loss updates the actor. The base stays frozen, and
fresh targets are constructed without overwriting stored transitions.
The target displacement satisfies
$\|\tilde r_i-\bar r_i\|_2\leq\delta_{\mathrm{tot}}$.
As noted in the problem statement, fitting these targets with a shared
network does not guarantee the same bound on realized policy changes.

The dead-zone field is proportional to the negative gradient of a hinged
radial penalty, $(\|r\|_2-\sqrt{D}\rho)_+^2$. Its geometry can therefore
also be placed in a conventional actor loss. Our hinge-loss control tests
this connection directly; target regression is an implementation of the
restoring principle, not its only possible form.

\subsection{Training and Variants}

We alternate online collection with critic and actor updates. Critic
training retains the baseline's ensemble temporal-difference updates,
$n$-step returns, and demonstration replay. These settings are matched
within each controlled actor comparison. The policy requires no action
likelihood, and inference uses the frozen generator plus the residual in
\eqref{eq:residual}.

The LIBERO recipe uses $K=16$, $\eta_0=0.05$,
$\eta_{\mathrm{anc}}=1$, $\rho=0.05$ rms, $\delta_Q=1$, and
$\delta_{\mathrm{tot}}=0.15$. We compare shared reward steps
$\eta_Q\in\{1,2\}$ across tasks. The robomimic recipe uses
$\eta_Q=0.5$ and $\rho=0$.
For reward-field ablations, we replace the action gradient with particle
transport~\cite{deng2026drifting} toward either the critic's top-$k$
candidates or candidates weighted by a softmax of standardized
$Q$-values. The clipping, restoration, and regression steps remain the same,
isolating how much the update benefits from the critic's gradient direction.
